\subsection{Regex Engines}

Using a transition table, it is straightforward 
to write a program to recognize strings in a regular 
language. If \texttt{T} is a matrix representing 
the transition table with state rows and character
columns, then it might look something like:
\begin{lstlisting}
    state = initial_state; //start state of FA
    while (true) {
        next_char = getc();
        if (next_char == EOF) break;

        next_state = T[state][next_char];

        if (next_state == ERROR) break;
        state = next_state;
    }
    if (is_final_state(state))
        //recognized a valid string
    else
        handle_error(next_char);
\end{lstlisting}

Up until now, we have only considered matching an entire string
to see if it is in a regular language.
If we want to match multiple tokens from a file,
we need to look ahead to see if the next character belongs to the
current token. If it does, we can continue. 
If it doesn't, the next character becomes part of 
the next token. We don't need to care about these 
details or tie-breaking or anything else when using 
ANTLR, as it takes care of everything behind the scenes.

After a program is tokenized, we need a parser to 
determine the structure of the program. 